\subsection{TENSORS AND INVARIANT FORMS}

Let G be a Lie group. We consider G as operating on itself by left (resp. right) translation. Let \(\lambda\) be a vector functor of class \(C^{\omega}\) for isomorphisms. Then \(\lambda(\mathrm{TG})\) is an analytic left (resp. right) vector G-bundle (§1, no. 8, Corollary to Proposition 16). The mapping \((g, u) \mapsto gu\) (resp. ug) of \(G \times \lambda(\mathrm{L}(G))\) onto \(\lambda(\mathrm{TG})\) is an isomorphism \(\phi\) (resp. \(\psi\) ) of vector G-bundles (§7, no. 8, Corollary 2 to Proposition 17). Every G-invariant section of \(\lambda(\mathrm{TG})\) is analytic and determined by its value at e (§1, no. 8, Corollary 1 to Proposition 17). Such a section is called left (resp. right) invariant. Let \(\sigma\) be a left invariant section of \(\lambda(\mathrm{TG})\) ; the transform \(\sigma'\) of \(\sigma\) under a right translation \(\delta(g)\) is defined by \(\sigma'(\delta(g)h) = \lambda(\mathrm{T}_{h}(\delta(g)))\sigma(h)\) for all \(h \in G\) ; it is also left invariant; it is also derived from \(\sigma\) by \(\gamma(g) \circ \delta(g) = \operatorname{Int}(g)\) and hence

\[
\sigma^ {\prime} (e) = \lambda (\mathrm{Ad} g). \sigma (e).\tag{26}
\]

Similarly, let \(\tau\) be a right invariant section of \(\lambda(\mathrm{TG})\) ; the transform \(\tau'\) of \(\tau\) under a left translation \(\gamma(g)\) is also right variant and

\[
\tau^ {\prime} (e) = \lambda (\mathrm{Ad} g). \tau (e).\tag{27}
\]

We now consider G × G as operating on G on the left by

\[
\left(\left(g, g ^ {\prime}\right), g ^ {\prime \prime}\right) \mapsto g g ^ {\prime \prime} g ^ {\prime - 1}.
\]

Then G is a left Lie homogeneous space of \(G \times G\) (§ 1, no. 6, Example). Hence \(\lambda(\mathrm{TG})\) is an analytic left vector \((\mathrm{G} \times \mathrm{G})\) -bundle. A section of \(\lambda(\mathrm{TG})\) is called biinvariant if it is invariant under the action of \(\mathrm{G} \times \mathrm{G}\) on \(\lambda(\mathrm{TG})\) , in other words if it is invariant under left and right translations. Let \(\lambda(\mathrm{L}(\mathrm{G}))_{0}\) be the set of elements of \(\lambda(\mathrm{L}(\mathrm{G}))\) invariant under \(\lambda(\mathrm{Ad}(\mathrm{G}))\) . For all \(u \in \lambda(\mathrm{L}(\mathrm{G}))_{0}\) , let \(\sigma_{u}\) be the mapping of G into \(\lambda(\mathrm{TG})\) defined by \(\sigma_{u}(g) = gu = ug\) . Then \(u \mapsto \sigma_{u}\) is a bijection of \(\lambda(\mathrm{L}(\mathrm{G}))_{0}\) onto the set of biinvariant sections of \(\lambda(\mathrm{TG})\) (§ 1, no. 8, Corollary 1 to Proposition 17).

PROPOSITION 50. Let G be a Lie Group (assumed to be finite-dimensional if K is of characteristic \textgreater0). Let E be the vector space of continuous alternating multilinear forms of degree k on \(T_{e}(G)\) . For all \(u \in E\) , let \(\omega^{u}\) be the differential form of degree k on G such that \((\omega^{u})_{g}\) is the multilinear form on \(T_{e}(G)\) derived from u by the translation \(h \mapsto gh\) (resp. \(h \mapsto hg\) ). Then \(\omega^{u}\) is analytic and left (resp. right) invariant on G. The mapping \(u \mapsto \omega^{u}\) is an isomorphism of E onto the vector space of left (resp. right) invariant differential forms of degree k on G.

This is a special case of what we have said above.

Let \(\mathbf{F}\) be a complete normable space. Proposition 50 remains true if differential forms on \(\mathbf{G}\) with values in \(\mathbf{K}\) are replaced by differential forms on \(\mathbf{G}\) with values in \(\mathbf{F}\) . For every continuous linear mapping \(u\) of \(\mathrm{T}_e(\mathrm{G})\) into \(\mathbf{F}\) , there exists a differential form \(\omega^u\) of degree 1 on \(\mathbf{G}\) , with values in \(\mathbf{F}\) , such that \((\omega^{u})_{g} = u \circ \mathrm{T}_{g}(\gamma(g)^{-1})\) . In particular, take \(\mathrm{F} = \mathrm{T}_e(\mathrm{G})\) and \(u = \mathrm{Id}_{\mathrm{T_e(G)}}\) . We then obtain the differential form \(\omega\) on G such that \(\omega_{g} = \mathrm{T}_{g}(\gamma (g^{-1}))\) ; this differential form is left invariant and analytic; it is called the left canonical differential form of G. \(\omega_{g}(t) = g^{-1}t\) for all \(t\in \mathrm{T}_g(\mathrm{G})\) .

If F is again an arbitrary complete normable space and \(u \in \mathcal{L}(\mathrm{T}_{e}(\mathrm{G}), \mathrm{F})\) , then \(\omega^{u} = u \circ \omega\) . In particular (taking F = K), the mapping \(v \mapsto v \circ \omega\) is a linear bijection of the dual of \(\mathrm{T}_{e}(\mathrm{G})\) onto the vector space of differential forms of degree 1 with values in K which are left invariant under G.

Similarly, the differential form \(\omega'\) on G such that \(\omega_g' = \mathrm{T}_g(\delta(g))\) is called the right canonical differential form of G. There are analogous properties to those of \(\omega\) , which we leave to the reader to state. The mapping \(g \mapsto g^{-1}\) of G onto G transforms \(\omega\) into \(\omega'\) .

